\begin{abstract}
Promptable video segmentation models such as SAM2 propagate an object mask through a video from a few expert prompts.
The propagated mask drifts under occlusion and appearance change, and repairing it requires further expert intervention.
Existing frame-selection methods decide where to place these corrections, but they fail to account for the trade-off between the resulting accuracy gain and the expert cost incurred in generating those corrections.
We formulate corrective prompting as a \textit{variable-cardinality learning-to-defer} problem: given the current propagation and a specified expert cost, our method decides whether to request assistance, how many corrections are justified, and which frames receive them.
A cardinality model estimates the number of frames to correct, and a multi-frame selection model selects the frames to be corrected conditioned on the frames already selected, so the requested prompts are complementary to one another.
Experiments with SAM2.1 on four datasets covering object transformation, camouflage, colonoscopy, and thyroid ultrasound show that our method achieves
the best quality--cost trade-off among learned methods.
Relative to the initial propagation, it improves mIoU by $5.6\%$--$47.8\%$.
Compared with the strongest
adapted baseline, Top-$k(x)$, it reduces corrective prompts by $3.1\%$--$9.2\%$ and system cost by $5.4\%$--$11.4\%$, while requesting and applying the selected corrections within a single correction round.
Code is released upon acceptance.
\end{abstract}
